% ==========================================================================
% Figure blocks for
%   Catalytic Ladders in a Federated Query Language: An Integration Report
%
% Each block below is self-contained. Paste one straight into
% the manuscript at the point its panel belongs; no \input and
% no \captionPanel* macro is required.
%
% Preamble requirements:
%   \usepackage{graphicx}
%   \usepackage{amsmath,amssymb}
%   \graphicspath{{figures/}}      % if the paths below are bare
%
% These captions contain no cross-references, so the blocks
% compile in any document that can show a figure.
% ==========================================================================


% --------------------------------------------------------------------------
\begin{figure*}[!htbp]
    \centering
    \includegraphics[width=\textwidth]{figures/panel_01.png}
    \caption{\textbf{The host's own suite, and the plans it ships, after the
    extension.}
    \textbf{(A)}~The pre-registered falsification condition. The host's suite
    comprises sixteen checks; sixteen pass and none fails, marked by the cross at
    zero. The suite is invoked as a subprocess and its own summary file is read, so
    the figure plots the host's report rather than our restatement of it. Had this
    bar been non-zero the integration would have failed and the failure would have
    been the result.
    \textbf{(B)}~The negative control, by fixture world. All eleven shipped plans
    that contain no ladder execute, across all three of the host's worlds. The
    world is chosen by the host's own selector; an earlier version of this control
    loaded a single world and reported three failures that were properties of the
    control rather than of the extension.
    \textbf{(C)}~Requests issued by each shipped plan. The plans still reach their
    sources and still spend budget: the control is not passing because nothing
    happens.
    \textbf{(D)}~Verdicts across the shipped plans, by step. The surface is
    concentrated on \textsf{answer} with a ridge at \textsf{starved}, which is the
    distribution the host's suite was written to produce and which the extension
    leaves in place.}
    \label{fig:panel1}
\end{figure*}


% --------------------------------------------------------------------------
\begin{figure*}[!htbp]
    \centering
    \includegraphics[width=\textwidth]{figures/panel_02.png}
    \caption{\textbf{The ladder occupies a cost tier the host already had.}
    \textbf{(A)}~Budget spent per step in a plan with two source accesses, one set
    operation and one ladder. The two source steps spend one request each; the
    union and the ladder spend nothing. The ladder's bar is not merely small, it is
    zero.
    \textbf{(B)}~Allocated against spent. Every non-source step is allocated the
    nominal unit the host's water-filling allocator assigns to all-or-nothing steps
    and then spends none of it. That asymmetry is the host's existing conservatism,
    visible here on the union step as much as on the ladder, and the ladder adopts
    it rather than introducing a third regime.
    \textbf{(C)}~Cumulative requests along the plan. The trace rises at each source
    step and is flat across the union and the ladder, so adding ladders to a plan
    cannot make it overrun its budget.
    \textbf{(D)}~Requests spent over the number of source steps and the number of
    ladder steps. The surface is a plane rising in one axis and constant in the
    other: cost depends on source accesses alone.}
    \label{fig:panel2}
\end{figure*}


% --------------------------------------------------------------------------
\begin{figure*}[!htbp]
    \centering
    \includegraphics[width=\textwidth]{figures/panel_03.png}
    \caption{\textbf{One composite, two verdicts, and a blame walk that stops where
    it should.}
    \textbf{(A)}~The same plan under two declared targets. Its rungs compose to
    $0.826\,75$ (horizontal line) in both cases; against a target of $0.70$ the
    step is answered and against $0.95$ it is refused. The verdict is a property of
    the declaration, not of the data, and the input step reports
    \textsf{answer} in both runs.
    \textbf{(B)}~The shortfall the diagnosis names: $0.95-0.826\,75=0.123\,25$,
    drawn as the part of the declared target the ladder does not reach. A refusal
    that named no quantity would be less useful than one that names this one.
    \textbf{(C)}~The blame walk. As implemented it returns the ladder alone; had
    the ladder named its input as the culprit it would have returned two steps and
    accused a step that answered correctly. The right-hand bar is the failure mode
    the empty predecessor field exists to prevent, and it is checked rather than
    assumed.
    \textbf{(D)}~The refusal surface: shortfall over composite and declared target.
    It is zero on and below the diagonal---where the ladder reaches its target---
    and rises above it. The two markers are the two runs of panel~A, one on the
    flat region and one on the slope.}
    \label{fig:panel3}
\end{figure*}


% --------------------------------------------------------------------------
\begin{figure*}[!htbp]
    \centering
    \includegraphics[width=\textwidth]{figures/panel_04.png}
    \caption{\textbf{The host's union bound is the additive relaxation of the
    ladder's law, and it goes vacuous as chains lengthen.}
    \textbf{(A)}~The two forms on identical retentions drawn from $[0.5,0.99]$,
    $4000$ trials per point. The multiplicative form stays near unity; the additive
    form falls, crosses zero between three and four stages, and enters the shaded
    region. A lower bound on a quantity known to lie in $[0,1]$ that has gone
    negative is true and says nothing.
    \textbf{(B)}~The fraction of trials on which the additive form is vacuous:
    zero at two stages, $0.515$ at four, $0.996$ at eight. Vacuity is the common
    case by four stages.
    \textbf{(C)}~The mean gap between the two forms, which grows monotonically from
    $0.44$ to $2.05$ across the range measured.
    \textbf{(D)}~The gap over chain length and retention. It vanishes at high
    retention and short chains---where the union bound is nearly tight---and grows
    in both directions away from that corner. The comparison is not a criticism:
    the additive form is stated as a bound and carries an injectivity hypothesis,
    and what is measured here is how much it gives up.}
    \label{fig:panel4}
\end{figure*}


% --------------------------------------------------------------------------
\begin{figure*}[!htbp]
    \centering
    \includegraphics[width=\textwidth]{figures/panel_05.png}
    \caption{\textbf{The frozen corpus: two public services with different schemas.}
    \textbf{(A)}~Sizes and coverage. Four hundred KEGG enzyme records, of which
    $372$ carry at least one reaction; $904$ Reactome reactions, of which $353$
    carry at least one catalyst. The catalyst count was zero on the first
    extraction and the discrepancy was found by querying the live service rather
    than by accepting a plausible zero.
    \textbf{(B)}~Reactions per pathway, ranked. The corpus supplies chains of
    genuinely differing length, the longest containing $58$ reactions, which is what
    makes the chain-length axis of panel~6 a property of data rather than of a
    parameter we chose.
    \textbf{(C)}~Reactions per enzyme record in KEGG, on a logarithmic count axis.
    The distribution is heavy-tailed: most records name few reactions and a few name
    many.
    \textbf{(D)}~Inputs, outputs and catalysts per Reactome reaction. The cloud
    sits mostly in the plane of zero catalysts, with a populated slab above it, and
    the two axes of participants are correlated but not equal---the schema-level
    difference from KEGG, which records substrates and products against an EC
    number instead.}
    \label{fig:panel5}
\end{figure*}


% --------------------------------------------------------------------------
\begin{figure*}[!htbp]
    \centering
    \includegraphics[width=\textwidth]{figures/panel_06.png}
    \caption{\textbf{Composition read against the corpus.}
    \textbf{(A)}~Composite power against the number of rungs at three uniform rung
    powers, with vertical lines at the six longest pathway lengths in the corpus.
    At a rung power of $0.05$ even the longest chain in the corpus falls short of
    unity, which is the saturation result placed against real chain lengths rather
    than against an arbitrary axis.
    \textbf{(B)}~Rungs required to reach three targets, as a function of rung
    power, on a logarithmic axis. The count diverges as the rung power falls, and
    the three curves separate sharply as the target approaches unity.
    \textbf{(C)}~Sensitivity against rung power for a single ladder of corpus
    length. Points lie exactly on $P/(1-\pi_j)$, which increases in $\pi_j$:
    control lies at the strongest rung, not the weakest.
    \textbf{(D)}~The saturation surface: composite power over chain length and
    rung power. It rises to unity in the far corner and stays low along the
    low-power edge however long the chain, which is the statement that a long
    sequence of weak steps does not substitute for a short sequence of strong
    ones.}
    \label{fig:panel6}
\end{figure*}

